% Dossier de compétences -- version orientée DevSecOps/Automatisation/Réseaux
\documentclass[10pt,a4paper]{article}
\usepackage[utf8]{inputenc}
\usepackage[T1]{fontenc}
\usepackage[default]{raleway}
\renewcommand*\familydefault{\sfdefault}
\usepackage{moresize}
\usepackage[a4paper,top=1.5cm,bottom=1.5cm,left=1.5cm,right=1.5cm]{geometry}
\usepackage{fancyhdr}
\pagestyle{fancy}
\setlength{\headheight}{-5pt}
\setlength{\parindent}{0mm}
\usepackage{multicol}
\usepackage{array}
\newcolumntype{x}[1]{>{\raggedleft\hspace{0pt}}p{#1}}
\usepackage{graphicx}
\usepackage{tikz}
\usetikzlibrary{shapes, backgrounds,mindmap, trees}
\usepackage{color}
% \definecolor{sectcol}{RGB}{0,150,255}
% \definecolor{bgcol}{RGB}{110,110,110}
% \definecolor{softcol}{RGB}{225,225,225}

% ADN colors
\definecolor{redADN}{RGB}{244,80,88} % official color palette
\definecolor{blueADN}{RGB}{102,142,212} % official color palette

\definecolor{sectcol}{RGB}{244,80,88}
\definecolor{bgcol}{RGB}{102,142,212}
\definecolor{softcol}{RGB}{225,225,225}

% small helper commands (adapted from template)
\renewcommand{\headrulewidth}{0pt}
\renewcommand{\footrulewidth}{0pt}
\renewcommand{\thepage}{}
\renewcommand{\thesection}{}
\newcommand{\cvsection}[1]{\begin{center}\large\textcolor{sectcol}{\textbf{#1}}\end{center}}
\newcommand{\cvevent}[5]{%
\begin{tabular*}{1\textwidth}{p{13.6cm}  x{3.9cm}}
	\textbf{#2} - \textcolor{bgcol}{#3} &   \vspace{2.5pt}\textcolor{sectcol}{#1}
\end{tabular*}
\vspace{-8pt}
\textcolor{softcol}{\hrule}
\vspace{6pt}
  $\cdot$ #4\\[3pt]
  $\cdot$ #5\\[6pt]
}

\begin{document}
\vspace{-8pt}
\begin{center}
  \HUGE \textcolor{sectcol}{\rule[-1mm]{1mm}{0.9cm}} \textsc{Dossier de compétences}\\[2pt]
  \small Spécialiste cybersécurité \& ingénierie DevSecOps
\end{center}
\vspace{6pt}

% META
\cvsection{Résumé}
Ingénieur en cybersécurité (M2 Cryptologie et Sécurité Informatique) avec expérience pratique en fuzzing embarqué, sécurité système embarqué, reverse engineering et automatisation. Orientation DevOps/DevSecOps : pipelines CI/CD, Infrastructure as Code, conteneurs et orchestration, automatisation des tests de sécurité, monitoring et conformité. Disponible immédiatement.

\vspace{4pt}
\noindent\textbf{Contact} : Paul Chauveau — Toulouse / La Rochelle — chauveaupaul17@gmail.com — +33 6 78 88 99 17 \\
\noindent\textbf{Langues} : Français (natif), Anglais B2/C1, Allemand A2

\cvsection{Compétences techniques (catégorisées)}
\begin{multicols}{2}
\textbf{DevSecOps \& Automation}\\
IaC : Terraform, Ansible (concepts, playbooks, roles), scripts d'automatisation (Bash, Python, PowerShell)\\
CI/CD : Git, GitLab CI/CD, GitOps (ArgoCD / Flux concepts), Jenkins, pipelines automatisés, tests intégrés\\
Containerisation \& Orchestration : Docker, images Dockerfile, Kubernetes (concepts, déploiement), Helm, Podman\\[6pt]

\textbf{Cloud \& Virtualisation}\\
Environnements : AWS / Azure (concepts), OpenStack / Proxmox (concepts)\\
Virtualisation \& émulation : QEMU, VirtualBox, KVM, Unicorn Engine (émulation CPU pour fuzzing)\\[6pt]

\textbf{Systèmes \& Réseaux}\\
GNU/Linux (Debian, Ubuntu, administration, kernel modules), systemd, SSH, Réseaux : VLAN, NAT, VPN, IPSec, routage, firewall (iptables/nftables)\\
SIEM/monitoring : SIEM (découverte), ELK stack, Prometheus, Grafana, alerting\\[6pt]

\textbf{Sécurité \& Cryptographie}\\
Protocoles : TLS, Kerberos, IPSec, PKI concepts, gestion de clés (NIST SP 800-57 awareness)\\
Cryptanalyses : DPA (side-channel), fault injection, attaques statistiques sur ECDSA, JavaCard security, cartes à puce\\
Hardening, audits, pentest basics (Metasploit), reverse engineering (Ghidra, GDB)\\[6pt]

\textbf{Testing \& Fuzzing}\\
Fuzzing embarqué : MCUboot fuzzing, AFL++, Unicorn-AFL, grammaires, custom mutators, AFL custom hooks, coverage analysis, sanitizers (ASAN/UBSAN)\\
Outils : afl, afl++ custom mutator, QEMU-based fuzzing, instrumentation, harness development\\[6pt]

\textbf{Langages \& Outils}\\
Python, Bash, PowerShell, C, C++, Java, Assembly (ARM / x86-64), Git, VSCode, Make/CMake, LaTeX, Markdown, Doxygen\\
\end{multicols}

\vspace{6pt}
\cvsection{Expériences sélectionnées orientées DevOps / DevSecOps}

\cvevent{Mars 2025 -- Août 2025}{STMicroelectronics — Stage de fin d'études}{Labège (Toulouse)}{Fuzzing du bootloader MCUboot et développement de grammars pour tests automatisés}{Sécurité système et logicielle des systèmes embarqués ARM. Mise en place de harness d'émulation (QEMU / Unicorn), intégration de campagnes de fuzzing automatisées; instrumentation pour coverage et rapports.}

\cvevent{Mai 2024 -- maintenant}{Solving Company — Mission / Consultant}{Bordeaux (freelance)}{Conformité ANSSI et RGPD pour traitements de données de santé; recommandations techniques}{Mise en place de bonnes pratiques DevSecOps, pipelines CI, automatisation des déploiements, gestion de secrets et conformité.}

\cvevent{Novembre 2023}{Crédit Agricole GIP — Immersion SOC (N3)}{Gradignan}{Parsing d'IOCs, analyse du SIEM Splunk}{Approche opérationnelle SOC, triage d'incidents, supervision.}

\cvevent{Mai 2021 -- Juil 2021}{Laboratoire XLIM UMR CNRS 7252 — Stage}{Poitiers}{Conception de travaux pratiques (C++)}{Développement d'exercices avancés, tests automatisés, packaging et documentation.}

\vspace{6pt}
\cvsection{Projets et réalisations (extraits)}
\begin{itemize}
  \item \textbf{MCUboot fuzzing (STMicro)} : création de grammars, harness d'émulation, campagnes fuzz automatisées, analyse de coverage et des plantages.
  \item \textbf{Custom mutator / AFL++ (travail universitaire / perso)} : définition de mutateurs structurés pour formats binaires, optimisation des queues, synchronisation multi-processus.
  \item \textbf{Hardening et conformité} : application des directives techniques ANSSI, audits RGPD pour données sensibles.
  \item \textbf{Reverse engineering} : usage de Ghidra et GDB pour l'analyse binaire, débogage et patching expérimental.
  \item \textbf{Systèmes embarqués} : modules noyau Linux (expérimentaux), assemblage et tests sur cibles ARM, outils arm-none-eabi.
\end{itemize}

\vspace{6pt}
\cvsection{Mots-clé (pour tri commercial / recherche de mission)}
\noindent\small
DevSecOps, CI/CD, GitLab CI, Jenkins, GitOps, Ansible, Terraform, IaC, Docker, Kubernetes, Helm, ArgoCD, Dockerfile, Prometheus, Grafana, ELK, SIEM, fuzzing, MCUboot, AFL++, Unicorn, Unicorn-AFL, QEMU, instrumentation, sanitizers, ASAN, UBSAN, coverage, custom mutator, embedded security, ARM, arm-none-eabi, kernel modules, Linux (Debian/Ubuntu), SSH, firewalls, VLAN, NAT, IPSec, TLS, Kerberos, PKI, JavaCard, DPA, fault injection, ECDSA, AES, reverse engineering, Ghidra, GDB, Python, Bash, C, C++, Assembly, LaTeX, Doxygen, RGPD, ANSSI, conformité.

\vspace{8pt}
\noindent\textbf{Disponibilité} : immédiate. \quad \textbf{Mobilité} : Bordeaux / Toulouse / France.

\null
\vspace*{\fill}
\end{document}
